% Wells Fargo Technology Program Intern 2027 | Off-campus technical interview
% Speak only what is on the printed resume in the interview folder (R-557386).
\documentclass[a4paper,9.5pt]{extarticle}
\usepackage[margin=0.85cm]{geometry}
\usepackage{enumitem}
\usepackage{titlesec}
\usepackage{xcolor}
\usepackage{tabularx}
\usepackage{array}
\usepackage[T1]{fontenc}

\pagestyle{plain}
\definecolor{ink}{HTML}{111827}
\definecolor{muted}{HTML}{4B5563}
\definecolor{accent}{HTML}{9F1239}
\definecolor{hot}{HTML}{9F1239}
\definecolor{ok}{HTML}{166534}
\color{ink}
\setlength{\parskip}{1.5pt}
\setlist[itemize]{leftmargin=1.1em, topsep=1pt, itemsep=1.4pt, parsep=0pt}
\titleformat{\section}{\bfseries\large\color{accent}}{}{0em}{}[\vspace{0.5pt}{\color{accent}\hrule height 0.55pt}]
\titlespacing{\section}{0pt}{8pt}{3pt}
\titleformat{\subsection}{\bfseries\normalsize}{}{0em}{}
\titlespacing{\subsection}{0pt}{5pt}{1.5pt}

\newcommand{\Q}[1]{\par\vspace{1.5pt}\noindent\textcolor{hot}{\textbf{Q:} #1}\par}
\newcommand{\Typ}[1]{\noindent\textcolor{muted}{\textbf{What people usually say:} #1}\par}
\newcommand{\Say}[1]{\noindent\textcolor{ok}{\textbf{You say:} #1}\par}

\begin{document}
\begin{center}
  {\LARGE\textbf{Wells Fargo Technology Program Intern 2027}}\\[2pt]
  {\small\color{hot}\textbf{Off-campus technical interview --- questions, the weak answer, and the answer you give}}\\[3pt]
  {\small Ojas Srivastava \textbar{} B.Tech AI, SVNIT Surat \textbar{} Resume freeze: requisition R-557386 \textbar{} Bengaluru / Hyderabad}\\[1pt]
  {\footnotesize\color{muted} Built 13 November 2026 from India 2025--2026 candidate reports plus the printed resume. Not a generic DSA dump. The earlier plan PDF is the background. This is the script.}
\end{center}

\section{0. How your loop actually works}
You applied \textbf{off campus}. No placement cell. A Wells Fargo caller has already scheduled a \textbf{technical interview} and asked whether you can be in the \textbf{Bangalore office on 5 and 6 December 2026} for something he called a hackathon. Those two dates are a Saturday and a Sunday.

\noindent\begin{tabularx}{\textwidth}{@{}>{\bfseries}p{3.3cm} X@{}}
What is public & India Technology Program Intern, 2028 graduating batch, summer 2027: resume, then an online assessment, then technical rounds, then HR. Bengaluru and Hyderabad. Desired skills on the posting: Java, C\#, SQL, Python, SDLC, spoken English. \\
What is yours & You are already at the live technical call. Treat the assessment as already passed or skipped. Do not spend tonight on a new HackerRank set. \\
What is not public & No Reddit, LinkedIn, Glassdoor, or careers page found that names a candidate hackathon on 5--6 December, or an off-campus Bangalore weekend as a required round. Internal intern hackathons exist. Those are for people already inside the internship. \\
\end{tabularx}

\textbf{The India shape, from people in your exact batch.} A takeUforward report for Technical Programme Intern, summer 2027, Bengaluru, on-campus, medium difficulty:
\begin{itemize}
  \item OA, 90 minutes, Intervue: 2 coding problems (an array ``minimum moves'' style, count palindromic substrings) plus about 10 MCQs on stocks and business charts. The writer said CGPA 8+ was the shortlist band. Yours is 9.20.
  \item Technical 1, 45--50 min: intro, AI vs ML, SQL vs PL/pgSQL, database vs schema, three SQL writes (second-highest salary, top 3 with a salary condition, highest-paid employee under manager E1), then Bubble Sort, stack, queue, linked list, and where you would use them.
  \item Technical 2, 25--30 min: almost entirely the resume. Challenges, how you fixed them, and, if the project touched money or prices, API limits and what the backend actually did.
  \item HR: family, academics, CPI, why Wells Fargo, what you would work on. That interviewer said team allocation can include GenAI and agentic work. It is not a promise.
\end{itemize}

\textbf{Other India reports that still match, and the ones that do not.}
Campus summer 2026 (GeeksforGeeks, September 2025): virtual technical on the project, kth largest, TCP vs UDP, process vs thread, deadlock, OOP; a second virtual technical that cross-examined the same project; a short HR. Nandini Sharma (Medium, September 2025): virtual technical with DP output, binary search, four OOP pillars, ArrayList vs LinkedList; then an \textbf{in-person senior manager} who read the resume line by line, including a hackathon win. LeetCode Discuss, Technology Summer Analyst Intern 2026: frequency count, compare two dates, HashMap vs array, joins, a SQL scenario, keys, indexing, then a second round that was the tech stack on the resume. August 2026 Reddit: OA had TSP and Dijkstra for some people; the interview asked tree-traversal \emph{types} and how hashing works, not a hard graph to code. Off-campus Technology Program Analyst, January 2026 (r/developersIndia): technical and techno-managerial the same day, then HR asked Bengaluru or Hyderabad and said about two weeks. People were still waiting after three.

\textbf{Do not import the US loop.} US SWE intern posts describe 30 minutes behavioural and one easy question on Zoom. Your caller is India, Bangalore office, requisition R-557386.

\textbf{Money, said once so you do not ask it.} A November 2025 on-campus intern posted Rs 90{,}000 a month, Hyderabad or Bengaluru. Aggregators for the 2027 posting print Rs 80{,}000--90{,}000. The official job text does not print a stipend. Ask nothing about pay in the technical round.

\textbf{What this resume allows you to claim.} Printed PDF only.
\begin{itemize}
  \item Nexus line is \textbf{DSA Mentor and Representative}. It does not say Chairperson. Do not say Chairperson.
  \item Community Hero does \textbf{not} say Global Top 20 on this PDF. Do not say Top 20.
  \item LeetCode line is Knight, rating 2048, \textbf{637+} solved. Do not say 711 or any other count.
  \item ``Neural Foundry'' is printed beside LogiFlow. If you cannot say what it is in one sentence, say ``that is the team name on the submission'' only if that is true, and stop. Do not invent a second product.
  \item Java, PostgreSQL, MongoDB, C, and C++ are in the skills block. Production bullets are Python, Node.js, MySQL, FastAPI, Redis, Firestore. Depth sentence first, then the concept.
\end{itemize}

\section{1. Locked intro (speak this, then stop)}
\textit{Good morning. I am Ojas Srivastava. I am a pre-final year B.Tech student in Artificial Intelligence at SVNIT Surat, with a CGPA of 9.20. I graduate in May 2028, and I am available from May to August 2027.}

\textit{Last summer I interned at IFFCO's Phulpur unit. I shipped production web services through the full lifecycle: REST APIs, MySQL with indexes and joins, JWT and input checks, Docker on Linux, and I had to explain the work to managers who do not write code.}

\textit{The project I want to walk through is LogiFlow, from Google Solution Challenge 2026, Global Top 106. I co-led the backend on Cloud Run. Corridor reads came back in 100 to 400 milliseconds through Redis, and I trained a delay model on 15{,}650 train-days.}

\textit{Wells Fargo is the same standard I already worked under. A service in front of a real user has to be correct, explainable, and safe when something fails. That is the work I want to do this summer.}

\vspace{2pt}
\noindent\textbf{Then stop.} Do not add Community Hero, AirHelp, the rating, Bengaluru, Hyderabad, the hackathon, or the stipend.\\
\textbf{If they ask ``go on'':} one sentence on Community Hero (solo civic-reporting API, Firestore, Gemini as an untrusted tagger) and stop again.

\section{2. Minute-by-minute: the technical call}
Reports put round one at 45--50 minutes. Yours may be shorter. Keep this order even if they jump.
\begin{enumerate}[leftmargin=1.3em, itemsep=2pt]
  \item \textbf{0--2 min.} Locked intro. Stop. Smile.
  \item \textbf{2--15 min.} They pick a noun you just said. Most likely LogiFlow, Redis, or the delay model. Have the 60-second LogiFlow answer (section 3). If they open the PDF themselves, IFFCO is the other door.
  \item \textbf{15--30 min.} Fundamentals they have actually asked this batch: AI vs ML (your branch), one SQL write, HashMap vs array, or OOP pillars. Sections 4 and 5.
  \item \textbf{30--45 min.} One Easy--Medium problem. Clarify, brute force, better approach, dry-run on a 5-element example, then code in \textbf{Python}. Complexity at the end. Section 6.
  \item \textbf{Last 3 min.} One question from section 8. No stipend. No ``is the hackathon eliminatory?'' unless they brought up 5--6 December.
\end{enumerate}

If a second interviewer joins the same day, that is the pattern from January 2026 off-campus and from Nandini's manager round: they re-read the resume and ask what \emph{you} owned. Do not repeat the intro at full length. Thirty seconds, then the project they did not already open.

\section{3. Project questions (this is where they decide)}
\subsection{LogiFlow --- you pointed them here}
\Q{Tell me about LogiFlow.}
\Typ{``It is an AI logistics app. We got Top 106. I used Python and Redis.''}
\Say{We helped an operator compare corridors on time, cost, and delay instead of getting one black-box route. Google Solution Challenge, Global Top 106. I co-led the backend on Cloud Run. Data on the resume: 580-plus corridors, 9{,}526 stations, 15{,}650 train-days. Hot reads sit in Redis so a repeat lookup is about 100 to 400 milliseconds. The delay model is one input to that ranking. pytest is 100 out of 100 business-rule checks. I did not write every frontend file.}

\Q{What was your code, and what was the team's?}
\Typ{``I was the lead of the whole product.''}
\Say{Co-led the backend. I own the read path, the Redis cache, the tests, and the delay-model training that the API exposes. A teammate owned parts of the UI and the submission. If I did not write a file, I say so.}

\Q{Why Redis? What if Redis is down?}
\Typ{``Redis is fast, so we used it.''}
\Say{Corridor metadata is read far more than it is written. Cache-aside: miss, then SQL, then set with a TTL. If Redis is down, FastAPI still answers from SQL, slower. Redis is not the source of truth. A stale map for a few minutes is acceptable. A stale balance or a stale permission is not. I would not cache those.}

\Q{MAE 22.7 minutes. Is that good?}
\Typ{``The model is very accurate.''}
\Say{On 15{,}650 train-days the mean absolute error is 22.7 minutes. That is the printed number. It means a typical prediction is about twenty minutes off a real delay, which is useful for ranking risk and useless as a promise to a passenger. I will not invent a baseline I did not run. If they want the model family and I can open the training code, I name it. If I cannot, I describe the target: predict delay, then feed it into the ranker. I do not call it deep learning.}

\Q{How does this relate to a bank?}
\Typ{``LogiFlow is basically Wells Fargo.''}
\Say{It is not a core-banking system. The transferable part is the discipline: a read path with a cache that can fail, tests on business rules, and a model that is an input, not the decision. At a bank I would keep the same split. The ledger is the source of truth. A cache and a model never override it.}

\subsection{IFFCO --- they will open this because it says production}
\Q{What did you actually ship?}
\Typ{``I did the full SDLC and built many APIs.''}
\Say{Phulpur unit, June to July 2025. Internal users, not farmers on the internet. I took one workflow from a requirement through design, a Node API, MySQL, a unit check, an integration check, and a Docker deploy on Linux. The resume says 10-plus APIs and 50-plus daily workflows. I will name \textbf{three endpoints and two workflows} when asked, with method, path, and who runs them. I will not hide behind the count.}

\Q{Show me the index.}
\Typ{``I added indexes and the query became fast.''}
\Say{Name the slow query, the column in the WHERE or the JOIN, the index, and what changed in time. If I do not remember the millisecond number, I say the before and after in plain language: full scan on that table versus a lookup. I do not invent a number.}

\Q{How does JWT work in your service?}
\Typ{``We used JWT for security.''}
\Say{The client sends a bearer token. The API checks the signature and the expiry before it touches MySQL. Missing or expired token: 401, no data. The token is not the user's password and it is not stored as the source of identity. Input that fails validation never reaches the query. That is the same rule I would want on a banking API.}

\Q{Explain one bug to a non-engineer.}
\Typ{A stack trace.}
\Say{Use the real triage. What the user saw, what was wrong in the data or the code, what you changed, how you checked it. No framework names until they ask.}

\subsection{Community Hero --- solo, so every word is yours}
\Q{What is it?}
\Typ{``An AI app with Gemini that solves civic problems.''}
\Say{A civic-report tool. I built it alone: React PWA, Express on Cloud Run, Firestore, Firebase Auth. A person drops a geo-tagged report. An admin triages it on a map. Gemini suggests a category. The resume lists geo-tagged reporting, admin workflows, map triage, Gemini categorization, SLA routing, community verification, and Open311 export, and it says eight features shipped. I list those, and I name the eighth from the hackathon brief only if I still remember it. I do not say Global Top 20. That rank is not on this resume.}

\Q{What if Gemini is wrong?}
\Typ{``We retry the model.''}
\Say{The model is untrusted. Empty or noisy text does not go to it. If the category is not one I know, I store the raw text and leave the report for an admin. A missed auto-tag is better than a wrong route. The API still returns a structured error, which is the line on the resume.}

\Q{Why Firestore, not MySQL?}
\Typ{``NoSQL is modern.''}
\Say{The report shape varies and I mostly fetch by id or by a geo query. Firestore matched that. IFFCO's workflows were tables, joins, and an index, so those stayed in MySQL. I do not pretend one store does both jobs.}

\subsection{AirHelp --- one bullet, many nouns}
\Q{Walk me through it.}
\Typ{``A* plus RAG plus WebSockets, 95 percent accurate.''}
\Say{Airport helper. FastAPI. Alerts go over a WebSocket because the client should not poll. The path on the map is A* over a graph: a node is a point a person can stand, an edge is a walkable link. The 1 to 3 seconds and the 50 to 100 users are what I measured, and I say how I measured them. 95 percent navigation accuracy is only as good as the test I describe. RAG: I retrieve a passage, and if it does not support the sentence, I do not show the sentence. NetworkX held the graph. I would not claim a production maps engine.}

\Q{Explain A* without a lecture.}
\Say{Best-first search. The score is cost so far plus a guess of the remaining cost. On my map the guess is straight-line distance. If the guess never overestimates, the first time I reach the goal I have a shortest path. BFS is what I would use if every corridor cost the same. I do not recite Dijkstra's code unless they ask for it.}

\subsection{Nouns they can point at}
\Q{Java. PostgreSQL. MongoDB. C++.}
\Say{Java: OOP coursework, I can explain a class, an interface, and an exception. I have not shipped a Java service. PostgreSQL: I know it as a SQL engine. My production SQL is MySQL. I have not written PL/pgSQL. MongoDB: say the real depth in one sentence, including ``I have not designed a production collection'' if that is the truth. C++: coursework or a small program, not the IFFCO service. C: same honesty.}

\Q{Neural Foundry. McKinsey Forward. RangRiti.}
\Say{Neural Foundry: one true sentence or ``I should not expand a name I cannot define.'' McKinsey.org Forward Fellow: a structured-communication fellowship, not a job at McKinsey. RangRiti, Web Wonders 2025: I was technical lead, and I name the page or feature I shipped.}

\Q{LeetCode Knight, 2048, 637+.}
\Say{Only if they ask. Then: ``Knight, 2048, six hundred thirty-seven plus problems. I would rather solve one problem here than talk about the rating.''}

\section{4. SQL and the questions this batch was given}
The summer-2027 Bengaluru round asked three writes. Practise these on paper until the pen does not stop.
\begin{enumerate}[leftmargin=1.3em, itemsep=2pt]
  \item \textbf{Second-highest salary.} \texttt{SELECT DISTINCT salary FROM employee ORDER BY salary DESC LIMIT 1 OFFSET 1;} Then say a tie: DISTINCT keeps one salary, and two people on that salary both exist in the table if they ask for names. For names, join back.
  \item \textbf{Top 3 with a condition.} Filter in WHERE first (salary above X, or department), then \texttt{ORDER BY salary DESC LIMIT 3}. HAVING is for groups, not for this.
  \item \textbf{Highest paid person under manager E1.} \texttt{WHERE manager\_id = 'E1' ORDER BY salary DESC LIMIT 1}. Say what you return if E1 has no reports: an empty result, not an error.
\end{enumerate}
\Q{SQL versus PL/pgSQL. Database versus schema.}
\Say{SQL is the query language. PL/pgSQL is PostgreSQL's language for functions and procedures with IF and loops. I write SQL. I have not shipped PL/pgSQL. A database is the whole catalog. A schema is a namespace inside it. PostgreSQL can hold many schemas in one database. In MySQL the words are often used for the same thing, which is why I say which engine I mean.}

\Q{Inner, left, index, keys.}
\Say{Inner keeps matches. Left keeps every left row and fills nulls where the right side misses. A primary key identifies a row and is not null. A foreign key points at a primary key. Unique forbids duplicates. An index speeds the lookup and slows the write. The IFFCO index is the example, not a definition.}

\section{5. AI versus ML, OOP, hashing, trees}
The same 2027 round opened on AI versus ML because the candidate's branch invited it. Yours will too.
\Q{What is the difference?}
\Typ{``AI is the future and ML is a subset, and my project uses AI.''}
\Say{AI is the aim: a system that takes a decision a person used to take. ML is one method: fit a model on data. LogiFlow's delay model is ML. Gemini on Community Hero is a model I call, not a model I trained. A rules check on JWT is neither. I do not call a cache ``AI''.}

\Q{OOP, from your code.}
\Say{Four pillars with one example each, from the FastAPI service or the Express API, not from a textbook animal. Encapsulation: the SQL stays behind a function the route calls. Abstraction: the caller asks for corridors and does not see the cache. Inheritance: only if you actually subclassed; otherwise say you composed services instead, and say why. Polymorphism: one handler shape, different behaviour by the real type, if you have it. Overloading is same name, different inputs, compile time in Java. Overriding is a child replacing a parent method. Abstract class can hold state. An interface is a contract. In Python the closest thing you use is a base class or a protocol. ArrayList versus LinkedList: index a lot, use the array; insert in the middle when you already hold the node, use the list.}

\Q{How does hashing work? Is a HashMap faster than an array?}
\Say{The key is turned into a bucket. Average lookup is constant. The worst case is a long chain in one bucket. An array scan is linear. The map wins when you look up many times, which is Two Sum and the frequency count. An array wins when the data is tiny or you need order by index. This is also the Redis story: a key to a value, not a scan of every corridor.}

\Q{Tree traversals, which the August 2026 interview asked as types, not as code.}
\Say{Inorder, preorder, postorder on a binary tree, and level order with a queue. Say when: inorder on a BST gives sorted order. Level order is what you want for ``nearest''. You do not need to code a red-black tree.}

\section{6. The coding problems they have actually used}
Speak, then type. Python. Brute force first if the optimal one is not immediate.
\begin{itemize}
  \item \textbf{Frequency of each element.} Map. If they want sorted output, say what you sort and the extra cost.
  \item \textbf{Compare two dates, dd/mm/yyyy.} Year, then month, then day. Say what you do with 31/02.
  \item \textbf{kth largest.} Min-heap of size k. If they ask for less memory, quickselect in words.
  \item \textbf{Longest common prefix.} Stop at the first mismatch. Empty list. First characters already differ.
  \item \textbf{Palindromic substrings} (this batch's OA). Expand around each centre, or the table. Count, do not list, unless they ask. Even and odd lengths.
  \item \textbf{Minimum moves so each value sits before the next consecutive} (this batch's OA, their wording was rough). Clarify with an example of four numbers before you code. Repeat the clarified statement back.
  \item \textbf{Linked-list cycle.} Slow and fast. One node. Cycle at the head.
  \item \textbf{Binary search.} Why the window shrinks. Missing target. Target at either end.
  \item \textbf{Best time to buy and sell stock,} one pass. A falling series returns 0. This is the closest ``finance'' coding question in the older intern reports. It is not options pricing. You do not have an options project on this resume. Do not bring Black--Scholes in.
  \item \textbf{Number of islands / a small BFS.} Only the idea, plus one traced grid. Your A* answer already shows you can separate BFS from a weighted search.
  \item \textbf{Bubble sort,} because this batch asked it as a concept. Say it is $O(n^2)$, stable, and you would not use it on the corridor table. Then say the sort you would use.
\end{itemize}

\section{7. One design question and one banking trade-off}
\Q{Design a token system for a branch.}
\Say{A queue. Each token stores a number and the time it was issued. One counter dequeues. Several counters: one queue, and a counter takes the head when it is free, so you do not have a customer stuck behind a slow counter. Say what you would log: who served whom, and when. That is the audit line a bank would want.}

\Q{The user wants it faster, and suggests dropping a security check.}
\Typ{``We can remove the check on the hot path.''}
\Say{The check stays. Auth and validation are not the cache. What I can speed up is what I already sped up on LogiFlow: cache data that is not sensitive, or fix the SQL, or move work off the request. I would measure the slow part before I touch the control. IFFCO's JWT check is the example of a control I would not cache my way around.}

\section{8. HR, location, and 5--6 December}
\Q{Why Wells Fargo?}
\Say{A large bank's software is used by people who cannot see the code. The program, as the posting describes it, starts with training and low-complexity work under a supervisor, inside compliance rules. That matches how I want to enter: I have shipped APIs, and I have not worked inside a bank's controls. I want the controls, not only the feature.}

\Q{Why should we hire you? Where will you work?}
\Say{I have already taken a service from a requirement to a deploy, and I can explain it to a non-engineer. Bengaluru and Hyderabad both work. Say that in one sentence. The resume header already says both. Do not negotiate a city in the technical round.}

\Q{Family, college, weakness, other processes.}
\Say{Short and factual. Prayagraj, school in Kota, college in Surat. One real weakness with a current fix, not a costume weakness. If they ask about other companies: this is the process I am in. Do not play offers against them.}

\Q{The caller asked about 5 and 6 December.}
\Say{In the reply, not as a speech in the technical round: yes or no based on whether you can actually be in Bangalore both days. Then ask, in writing, whether it happens only after this interview, whether it is solo or a team, what to carry, the reporting time and office, and whether travel is covered. Build nothing until that mail describes a format. If the interviewer raises it, repeat that you have replied and you are waiting for the written format. Do not pitch a hackathon idea.}

\textbf{Questions you ask them, if they offer the minute.}
\begin{itemize}
  \item What does an intern in this program ship in the first month, as opposed to the training?
  \item How is a team split between the Bengaluru and Hyderabad offices?
  \item What would you look at in week four to decide the internship is going well?
\end{itemize}

\section{9. Weak versus strong}
\noindent\begin{tabularx}{\textwidth}{@{}>{\bfseries}p{3.1cm} >{\color{muted}}X >{\color{ok}}X@{}}
Topic & Weak & Strong \\
Intro & Three minutes, every project, the rating, both cities & Locked intro, stop on LogiFlow \\
LogiFlow & ``AI app, I did everything, MAE is very good'' & Co-led backend, Redis can fail, 22.7 min is a rank signal \\
IFFCO & ``10 APIs, 50 workflows, full SDLC'' & Three endpoints, one index, one bug, one non-engineer explanation \\
Gemini & ``We used AI'' & Untrusted tagger, fallback to the admin \\
Java / Postgres & Pretend you shipped them & Coursework or familiarity, production stack named \\
DSA & Silent, then the optimal code & Clarify, brute, dry-run, Python \\
Don't know & Invent PL/pgSQL or a PPO rate & ``I have not shipped that; the closest thing I have is\ldots'' \\
December & Start building a product tonight & Reply sent, format unknown, interview first \\
\end{tabularx}

\section{10. Day-of rules}
\begin{itemize}
  \item The printed resume in the room is the only source of truth. Chairperson, Global Top 20, and any LeetCode count other than 637-plus are not on it.
  \item Every noun you add becomes a question: Redis, pytest, JWT, Gemini, A*, RAG, Neural Foundry.
  \item Code in Python. If they ask for Java syntax, say you will write the logic in Python and name the Java type (a \texttt{Deque}) if you know it. Do not freeze on syntax you do not use.
  \item Partial spoken brute force beats a silent clever answer.
  \item No stipend, no PPO percentage, no ``how hard is the hackathon'' in this round.
  \item Third-party recordings are prohibited unless Wells Fargo authorises them. Do not record the call.
  \item Official mail is \texttt{@wellsfargo.com}. They do not ask for payment or bank details.
  \item After you hang up, write the questions they actually asked. Send a document only if they asked for one. Keep other applications moving. Off-campus HR has taken weeks after a clear.
\end{itemize}

\section{11. Sources (so this is not invented)}
\begin{itemize}
  \item Your printed resume, requisition R-557386, file dated 13 November 2026 in the interview folder. Header: Bengaluru / Hyderabad. Coursework includes OOP in Java, DBMS, OS, networks. LeetCode 637-plus. Nexus: DSA Mentor and Representative. Community Hero has no Top-20 line.
  \item takeUforward, anonymous, about September 2026: Technical Programme Intern, summer 2027, Bengaluru, on-campus. OA on Intervue. R1 AI/ML, SQL vs PL/pgSQL, database vs schema, three SQL writes, Bubble Sort and linear structures. R2 resume, including a finance backend. HR mentioned GenAI and agentic work as team-dependent.
  \item Job text repeated by campus trackers for the 2027 Technology Program Intern: graduation March 2028--August 2028, internship May--August 2027, desired Java, C\#, SQL, Python, SDLC. Stipend printed by those trackers as Rs 80{,}000--90{,}000. The Workday text itself does not state pay.
  \item r/Btechtards, 23 November 2025: on-campus intern, Rs 90{,}000 a month, Hyderabad or Bengaluru, ``DSA was the focus.'' Comments asked for projects and did not get a question list. Treat the stipend as one candidate's post.
  \item r/developersIndia, March 2026: off-campus Technology Program Analyst. Technical plus techno-managerial on 29 January, HR on 2 February said both rounds were clear, asked Bengaluru or Hyderabad, ``about two weeks.'' A comment the next day: still waiting at three weeks.
  \item r/Btechtards, August 2026: OA included TSP and Dijkstra for two people. Interview: tree-traversal types and how hashing works.
  \item r/Btechtards, August 2025: seniors said number of islands; another comment: graphs and DP; another: Java and Python, C++ not allowed on that OA, 2 DSA and 1 SQL.
  \item GeeksforGeeks, Wells Fargo on-campus, September 2025 (updated January 2026): kth largest, routing, TCP vs UDP, process vs thread, deadlock, OOP, then HR.
  \item Nandini Sharma, Medium, 30 September 2025: virtual technical, then in-person senior manager, resume read line by line, hackathon project included.
  \item LeetCode Discuss, Abhishek Godara, May 2025, Technology Summer Analyst Intern 2026: longest common prefix, Euler path on the OA, then frequency count, date compare, HashMap vs array, joins, indexing.
  \item Srishtik Dutta, LinkedIn, February 2025, NIT Durgapur intern drive: longest common substring, stock profit, token system for a bank.
  \item Aditya Kumar, LinkedIn, 2026: applied on the careers portal with no referral. Rounds covered DSA, OOP, OS, DBMS, networks, projects.
  \item Anubhav Mishra, LinkedIn, 2026: off-campus Technology Program Analyst. OA, then two technical interviews, DSA, fundamentals, system design, business cases.
  \item r/LeetcodeDesi, July 2026: OA with SQL and coding, technical round on moderate DSA plus resume, techno-managerial, HR. ``Not too advanced compared with product companies.''
  \item Wells Fargo, Our Hiring Process: format varies by business. Phone, group, or one-to-one.
  \item Workday posting for the 2027 Technology Internship, Bengaluru: training, low-complexity work, supervision, compliance. Third-party recordings prohibited unless authorised. You represent your own experience.
  \item Searched November 2026 for a public 5--6 December Bangalore candidate hackathon. Nothing found. Do not treat the weekend as a known format.
\end{itemize}

\vspace{2mm}
\noindent{\footnotesize\color{hot}\textbf{Rule:} If it is not on the printed resume and you cannot dry-run it, do not say it. Clarify, then approach, then code, then a dry-run, then complexity. Never silent.}
\end{document}
